\documentclass[10pt]{article}

\usepackage[a4paper,margin=0.85in]{geometry}
\usepackage{amsmath,amssymb}
\usepackage{booktabs}
\usepackage{longtable}
\usepackage{array}
\usepackage{tabularx}
\usepackage{graphicx}
\usepackage{caption}
\usepackage{float}
\usepackage{pdflscape}
\usepackage{xcolor}
\usepackage[colorlinks=true,linkcolor=blue,citecolor=blue,urlcolor=blue]{hyperref}
\usepackage{titlesec}
\usepackage{fancyhdr}
\usepackage{parskip}

\graphicspath{{../data/results/run_20260526_v11_dct1_phospho_mediation/figures/v11/}{figures_v4/}{figures_osd462/}}
\newcommand{\gene}[1]{\textit{#1}}
\newcommand{\code}[1]{\texttt{\detokenize{#1}}}

\pagestyle{fancy}
\fancyhf{}
\fancyhead[L]{\small Curated results compendium -- Manuscript v11}
\fancyhead[R]{\small S\thepage}
\renewcommand{\headrulewidth}{0.4pt}
\renewcommand{\thepage}{S\arabic{page}}
\renewcommand{\thesection}{S\arabic{section}}
\renewcommand{\thetable}{S\arabic{table}}
\renewcommand{\thefigure}{S\arabic{figure}}
\makeatletter
\renewcommand{\@pnumwidth}{3.0em}
\renewcommand{\@tocrmarg}{4.0em}
\makeatother

\title{\textbf{Curated Results, Tables, and Figure Compendium}\\[0.3em]
\large Companion to Manuscript v11 -- distal-nephron subtype-prior phosphosite enrichment in mouse spaceflight kidney}
\author{Ibrahim Shahid}
\date{Draft v11 results compendium, May 2026}

\begin{document}
\maketitle
\thispagestyle{fancy}

\noindent
This document collects the curated results, manuscript tables, figure gallery, and artifact index supporting Manuscript v11. It is the technical supplement for the streamlined manuscript: cross-cohort RNA recurrence, matched OSD-462 RNA/protein/phosphoprotein mismatch, distal-nephron DCT-subtype-prior phosphosite enrichment, parent-protein and compartment-score sensitivity models, exploratory covariance decomposition, external spatial-reference context, and the public perturbation-reference pass (low-K DCT, parent-protein-normalized phosphosite effects, IRI DCT-transport spatial context, dDAVP/mpkDCT, KLHL3/CUL3 turnover). The main manuscript carries the story; this compendium carries the extra tables, secondary analyses, and run artifacts.

\tableofcontents
\clearpage

% =====================================================================
\section{Run provenance and executable entrypoint}

\begin{table}[H]
\centering\small
\caption{\textbf{V11 run provenance.}}
\label{tab:provenance}
\begin{tabularx}{\textwidth}{>{\bfseries}p{0.26\textwidth}X}
\toprule
Run root & \path{data/results/run_20260526_v11_dct1_phospho_mediation/} \\
Plan & \path{docs/v11_execution_research_plan.md} \\
Execution report & \path{docs/v11_execution_results.md} \\
Plan checksum & \code{d8b7465e09257dd4444fc2abb969a02066f2d360e4fbc04056791a60c9755e3f} \\
Pipeline entrypoint & \code{python -m src.run_all_phases --v11-only --run-id <run_id>} \\
Canonical v11 modules & \path{src/v11/build_gse228367_dct_prior.R}; \path{src/v11/core_analysis.py}; \path{src/v11/h2_composition_aware_phospho.py}; \path{src/v11/spatial_reference_projection.py}; \path{src/v11/publication_figures.py} \\
Public perturbation-reference modules & \path{src/v11/perturbation_gse228367_lowk.py}; \path{src/v11/h2_occupancy_normalized_phospho.py}; \path{src/v11/spatial_dct_transport_check.py}; \path{src/v11/perturbation_triage_summary.py} \\
\bottomrule
\end{tabularx}
\end{table}

\begin{table}[H]
\centering\small
\caption{\textbf{External datasets added or reused for v11.}}
\label{tab:datasets}
\begin{tabularx}{\textwidth}{p{0.20\textwidth}p{0.24\textwidth}X}
\toprule
Layer & Dataset & V11 role \\
\midrule
Spaceflight RNA recurrence & RRRM-2 / OSD-771, OSD-513, OSD-462 & Defines and tests the recurrent matrix/endothelial-high and DCT-low RNA state. \\
Spaceflight multi-omics & OSD-462 / RR-10 & Whole-kidney matched RNA, protein, and phosphoprotein anchor. \\
DCT subtype prior & GSE228367 & Native kidney DCT1/DCT2 snRNA-seq reference used to score parent-gene DCT1-high and DCT2-leaning priors. \\
DCT-lineage phosphoproteome & PXD001729 & mpkDCT dDAVP phosphoproteomic reference; used for DCT-lineage plausibility and negative anti-alignment boundary. \\
Spatial reference & GSE269622 & Whole-transcriptome Visium IRI timecourse used for pathway-vector projection. \\
Spatial annotation context & GSE269719 & Targeted 299-gene Xenium object used for cell-type and cellular-neighborhood inventory. \\
\bottomrule
\end{tabularx}
\end{table}

% =====================================================================
\section{Executive verdict and claim labels}

\begin{table}[H]
\centering\small
\caption{\textbf{Main result verdicts.}}
\label{tab:verdicts}
\begin{tabularx}{\textwidth}{p{0.17\textwidth}Xp{0.25\textwidth}}
\toprule
Hypothesis & Question & Verdict \\
\midrule
RNA recurrence & Does a recurrent matrix/endothelial-high and DCT-low RNA remodeling state recur across cohorts and in OSD-462 RNA? & Supported; remains the backbone. \\
DCT-subtype-prior enrichment & Are OSD-462 flight-suppressed phosphosites enriched among distal-nephron subtype-prior parent-gene extremes? & Supported for DCT1 top-decile and DCT2-bottom-decile bins; DCT1 is not exclusive. \\
Sensitivity models & Does the enrichment remain after parent-protein, parent-gene, observability, and compartment-score sensitivity checks? & Top-decile DCT1 subset survives many checks but attenuates in matched parent-gene permutation; DCT2-bottom is stronger in several parent-gene-aware summaries. \\
Exploratory covariance decomposition & Are remodeling scores statistically related to NCC/SPAK phospho suppression? & Endothelial and matrix/endothelial covariance summaries are directionally consistent; temporal order remains experimental. \\
Spatial supplement & Which injury/repair spatial niches resemble the bulk spaceflight RNA vector? & Contextualized to repair-stage, DCT-adjacent/injured-tubule and fibro-inflammatory/endothelial niches. \\
Public perturbation references & Do public perturbation references (low-K DCT, parent-protein normalization, IRI DCT-transport, dDAVP, KLHL3/CUL3) identify the upstream WNK-suppressing mechanism? & No single upstream signal identified; parent-protein-normalized DCT-prior enrichment is a robustness upgrade; IRI Visium DCT-adjacent transport drop is a testable spatial prediction. \\
\bottomrule
\end{tabularx}
\end{table}

\begin{table}[H]
\centering\small
\caption{\textbf{Preferred and avoided manuscript language.}}
\label{tab:language}
\begin{tabularx}{\textwidth}{>{\bfseries}p{0.24\textwidth}XX}
\toprule
Topic & Preferred wording & Avoid \\
\midrule
Overall frame & Distal-nephron subtype-prior phosphosite parent-gene enrichment within recurrent kidney remodeling. & DCT1-anchored mechanism without qualification. \\
DCT-prior enrichment & DCT1 top-decile and DCT2-bottom-decile enrichment of flight-suppressed whole-kidney phosphosites. & DCT1-specific phosphoproteomics or DCT1 exclusivity. \\
Composition models & Bulk sensitivity adjustment for parent protein abundance and estimated compartments. & Fully deconfounded or deconvolved cell-of-origin analysis. \\
Covariance decomposition & Remodeling scores covary with NCC regulatory phosphorylation after conditioning choices. & Endothelial remodeling causes NCC dephosphorylation. \\
Spatial & External IRI spatial reference contextualization. & Spatial validation of the spaceflight lesion. \\
\bottomrule
\end{tabularx}
\end{table}

% =====================================================================
\section{RNA recurrence and RNA-protein mismatch}

\begin{table}[H]
\centering\small
\caption{\textbf{Core RNA recurrence and multi-omic quantities from the v11 run.}}
\label{tab:h1}
\begin{tabularx}{\textwidth}{p{0.30\textwidth}rX}
\toprule
Component & Value & Interpretation \\
\midrule
OSD-462 RNA recurrence & 0.869 & OSD-462 recurs the RRRM-2 matrix-high / DCT-low direction. \\
OSD-462 protein abundance & null & RNA remodeling does not propagate cleanly as protein abundance. \\
NCC regulatory phospho score & -0.832 & NCC regulatory phosphosites are suppressed while total NCC protein is flat. \\
Targeted KSEA SPAK/OSR1 & -6.31 & Inferred kinase-output down; three quantified curated substrates. \\
Targeted KSEA WNK & -4.12 & Inferred WNK-output down; three quantified curated substrates, interpreted as pathway coherence. \\
Endothelial score vs NCC phospho & -0.762 & Strongest compartment anti-correlation with NCC regulatory phosphorylation. \\
\bottomrule
\end{tabularx}
\end{table}

\noindent
Interpretation: recurrent RNA remodeling is strong, targeted protein-abundance concordance is null, and the transporter lesion resolves at the phospho/activity layer.

% =====================================================================
\section{GSE228367 DCT subtype prior}

\begin{table}[H]
\centering\small
\caption{\textbf{GSE228367 DCT reference inventory.}}
\label{tab:dctinventory}
\begin{tabularx}{0.78\textwidth}{>{\bfseries}p{0.36\textwidth}r}
\toprule
DCT1 nuclei & 18,881 \\
DCT2 nuclei & 6,274 \\
Normal-potassium replicates per subtype & 3 \\
Strict FDR DCT1 core genes & 2 \\
Strict FDR DCT2 core genes & 0 \\
\bottomrule
\end{tabularx}
\end{table}

\begin{table}[H]
\centering\small
\caption{\textbf{Selected marker sanity checks in the DCT1/DCT2 reference.}}
\label{tab:dctmarkers}
\begin{tabular}{lrrl}
\toprule
Gene & DCT1 mean & DCT2 mean & Read \\
\midrule
\gene{Slc12a3} & 5.14 & 4.40 & present in both, higher in DCT1 \\
\gene{Pvalb} & 0.279 & 0.024 & DCT1-high \\
\gene{Trpm6} & 2.40 & 1.90 & DCT1-high \\
\gene{Klhl3} & 2.84 & 2.65 & mild DCT1-high \\
\gene{Trpv5} & 0.006 & 0.711 & DCT2-high \\
\gene{Calb1} & 1.01 & 2.97 & DCT2-high \\
\gene{Stk39} & 1.98 & 1.54 & DCT1-leaning \\
\gene{Wnk4} & 1.72 & 1.40 & DCT1-leaning \\
\bottomrule
\end{tabular}
\end{table}

\noindent
Interpretation: the strict FDR marker sets are too sparse for binary enrichment. V11 therefore uses continuous DCT1/DCT2 subtype-prior scores and percentile bins. Leave-one-replicate-out mean-difference priors were stable for the DCT1 top decile and DCT2 bottom decile (Jaccard approximately 0.85--0.86), while log-ratio and rank-average score definitions changed bin membership more strongly.

% =====================================================================
\section{Distal-nephron subtype-prior phosphosite enrichment}

\begin{figure}[H]
\centering
\includegraphics[width=\textwidth]{v11_dct1_parent_gene_enrichment.pdf}
\caption{\textbf{Distal-nephron subtype-prior parent-gene enrichment.} DCT1 top-decile, DCT2-bottom-decile, and DCT1 top-quartile reference bins are tested across primary and sensitivity definitions. The safest claim is enrichment in distal-nephron subtype-prior extremes rather than DCT1 exclusivity.}
\label{fig:h2primary}
\end{figure}

\begin{table}[H]
\centering\small
\caption{\textbf{Primary DCT1 top-decile enrichment tests within the DCT-subtype-prior analysis.}}
\label{tab:h2primary}
\begin{tabularx}{\textwidth}{p{0.34\textwidth}p{0.22\textwidth}rlX}
\toprule
Test & Suppressed set & Result & q & Read \\
\midrule
Matched-null mean DCT1 score & p<0.05 suppressed sites & -- & 0.065 & weak continuous support only \\
Fisher DCT1 top quartile & p<0.05 suppressed sites & OR 1.14 & 0.0061 & supported but weaker than top decile \\
Fisher DCT1 top decile & p<0.05 suppressed sites & OR 1.51 & $1.33\times10^{-11}$ & supported at row level \\
DCT1 top decile, anchor genes excluded & p<0.05 suppressed sites & OR 1.49 & $2.95\times10^{-11}$ & survives anchor exclusion \\
DCT1 top decile, NCC sites excluded & p<0.05 suppressed sites & OR 1.49 & $2.95\times10^{-11}$ & survives NCC-site exclusion \\
Composite/multi-position sites excluded & p<0.05 suppressed sites & OR 1.38 & $2.0\times10^{-6}$ & removes composite rows; not a parent-gene dependence fix \\
One row per parent gene & p<0.05 suppressed sites & OR 1.49 & $6.59\times10^{-4}$ & survives parent-gene representative-row sensitivity \\
Single-position row per parent gene & p<0.05 suppressed sites & OR 1.48 & $9.05\times10^{-4}$ & stricter representative-row sensitivity passes \\
Parent-gene Fisher & parent has at least one p<0.05 suppressed site & OR 1.52 & $1.90\times10^{-4}$ & parent-gene-level confirmation \\
Parent-gene logistic & parent has at least one p<0.05 suppressed site & OR 1.14 & 0.424 & attenuated after site-count, peptide-count, abundance, and missingness covariates \\
Matched parent-gene permutation & parent has at least one p<0.05 suppressed site & OR 1.51 & $p_{\mathrm{emp}}=0.107$ & attenuated after observability matching \\
Cluster bootstrap & row-level p<0.05 suppressed sites & OR 1.51 & 1.22--1.81 & parent-gene cluster-bootstrap interval excludes one \\
\bottomrule
\end{tabularx}
\end{table}

\begin{table}[H]
\centering\small
\caption{\textbf{DCT2-bottom-decile comparator across the same sensitivity ladder.} The DCT2-bottom-decile bin contains the 10\% of detected parent genes with the lowest DCT1 enrichment score (most DCT2-leaning). Both DCT1-top-decile and DCT2-bottom-decile use the same one-sided Fisher exact alternative. In several parent-gene-aware summaries the DCT2-bottom-decile odds ratio is larger than the DCT1-top-decile.}
\label{tab:dct2comparesupp}
\resizebox{\textwidth}{!}{%
\begin{tabular}{lrrrrl}
\toprule
Analysis & DCT1 OR & DCT1 q & DCT2 OR & DCT2 q & Read \\
\midrule
Primary $p<0.05$ & 1.51 & $1.33\times10^{-11}$ & 1.29 & $1.30\times10^{-5}$ & DCT1 stronger \\
Anchor genes excluded & 1.49 & $2.95\times10^{-11}$ & 1.29 & $1.30\times10^{-5}$ & DCT1 stronger \\
NCC sites excluded & 1.49 & $2.95\times10^{-11}$ & 1.30 & $1.30\times10^{-5}$ & DCT1 stronger \\
Composite rows excluded & 1.38 & $2.0\times10^{-6}$ & 1.33 & $1.30\times10^{-5}$ & similar \\
One row per parent gene & 1.49 & $6.59\times10^{-4}$ & \textbf{1.68} & $\mathbf{4.62\times10^{-6}}$ & DCT2 larger \\
Single-position row per gene & 1.48 & $9.05\times10^{-4}$ & \textbf{1.52} & $\mathbf{2.91\times10^{-4}}$ & both pass \\
Parent-gene Fisher & 1.52 & $1.90\times10^{-4}$ & \textbf{1.82} & $\mathbf{3.29\times10^{-8}}$ & DCT2 larger \\
Matched parent-gene permutation & 1.51 & 0.107 & \textbf{1.70} & \textbf{0.017} & DCT1 attenuates \\
\bottomrule
\end{tabular}
}
\end{table}

\noindent
Interpretation: suppressed whole-kidney OSD-462 phosphosites are enriched among parent genes in distal-nephron subtype-prior extremes. The DCT1 top decile is biologically motivated by NCC/SPAK/WNK anchors, while the DCT2-bottom-decile comparator is also enriched and is stronger in several parent-gene-aware summaries. The supported claim is therefore distal-nephron subtype-prior enrichment rather than DCT1 exclusivity. Direct cell-of-origin assignment would require DCT-enriched or spatially resolved spaceflight phosphoproteomics.

% =====================================================================
\section{Parent-protein and composition-score sensitivity}

\begin{figure}[H]
\centering
\includegraphics[width=\textwidth]{v11_parent_protein_composition_sensitivity.pdf}
\caption{\textbf{Parent-protein and composition-score sensitivity.} Panel A shows top-decile and top-quartile enrichment across the adjustment ladder. Panel B shows that continuous DCT1-reference and flight-by-DCT1 interaction models remain weak.}
\label{fig:h2composition}
\end{figure}

\begin{table}[H]
\centering\small
\caption{\textbf{DCT1 top-decile adjusted-suppression enrichment across the composition-aware model ladder.}}
\label{tab:h2ladder}
\begin{tabularx}{\textwidth}{p{0.12\textwidth}XrrX}
\toprule
Model & Adjustment & OR & q & Read \\
\midrule
M0 & Raw flight effect & 1.40 & 6.35e-7 & raw DCT1 top-decile subset enrichment \\
M1 & Parent protein abundance & 1.55 & 2.14e-10 & not eliminated by parent protein abundance adjustment \\
M2 & DCT score & 1.36 & 2.17e-5 & survives estimated DCT abundance adjustment \\
M3 & Endothelial + stromal scores & 1.24 & 0.00545 & attenuated, still positive \\
M4 & Parent protein + DCT + endothelial + stromal & 1.30 & 0.00158 & full model remains positive \\
M5 & Parent protein + composition PC1 & 1.38 & 6.53e-5 & PC-based composition adjustment passes \\
\bottomrule
\end{tabularx}
\end{table}

\begin{table}[H]
\centering\small
\caption{\textbf{Continuous DCT1-reference models do not pass.}}
\label{tab:h2continuous}
\begin{tabularx}{\textwidth}{p{0.30\textwidth}p{0.20\textwidth}rrX}
\toprule
Test & Key term & Coefficient & q & Read \\
\midrule
Two-stage M0 DCT1-only & DCT1 prior & -0.00110 & 0.877 & weak negative, not significant \\
Two-stage M4 full & DCT1 prior & +0.00274 & 0.877 & no continuous suppression gradient \\
One-shot LM0 site-fixed & flight x DCT1 prior & -0.00145 & 0.665 & direction negative, weak \\
One-shot LM4 full site-fixed & flight x DCT1 prior & -0.00200 & 0.665 & direction negative, weak \\
\bottomrule
\end{tabularx}
\end{table}

\noindent
Interpretation: the robust result is top-decile subset enrichment after parent-protein and compartment-score adjustment, not a broad continuous DCT1-gradient effect. The adjusted models redefine the suppressed set on each residual scale (M0: 2,430 nominally suppressed single-position sites; M4: 1,390). A fixed-M0-set sensitivity showed that 99.3\% of M0 suppressed sites with M4 covariates remained negative in the full model (mean adjusted effect $-0.295$), including 99.7\% of DCT1 top-decile sites (mean $-0.303$). Sample-level covariate diagnostics showed endothelial--stromal correlation $r=0.86$ and stromal-score VIF 6.59, so these are bulk sensitivity models rather than deconfounding.

% =====================================================================
\section{PXD001729 and KLHL3/CUL3 checks}

The PXD001729 dDAVP/mpkDCT anti-alignment check and the KLHL3/CUL3/WNK turnover coverage check are presented under the public-perturbation-reference section (\S\ref{sec:perturbation} below), as Branches~4 and~5 of the triangulation pass. See Tables~\ref{tab:pxdshared} and~\ref{tab:klhlcoverage} for the shared-site cosine and the OSD-462/PXD coverage numbers, and the corresponding interpretation paragraphs.

% =====================================================================
\section{Exploratory covariance decomposition}

\begin{figure}[H]
\centering
\includegraphics[width=0.92\textwidth]{v11_exploratory_covariance_decomposition.pdf}
\caption{\textbf{Exploratory remodeling/transporter covariance decomposition.} Endothelial and matrix/endothelial composite summaries are negative with intervals excluding zero under the implemented conditioning choices. Cross-sectional bulk data prevent causal interpretation.}
\label{fig:h3}
\end{figure}

\begin{table}[H]
\centering\small
\caption{\textbf{Covariance-decomposition summaries from approximate Bayesian OLS models.}}
\label{tab:h3}
\begin{tabularx}{\textwidth}{p{0.30\textwidth}rrrX}
\toprule
Intermediate score & Summary median & 95\% low & 95\% high & Read \\
\midrule
Endothelial & -0.432 & -1.019 & -0.036 & negative remodeling-phosphorylation covariance summary \\
Stromal/fibroblast & -0.295 & -0.843 & 0.025 & suggestive, interval crosses zero \\
DCT identity & -0.187 & -0.658 & 0.075 & composition-linked, interval crosses zero \\
Matrix/endothelial composite & -0.408 & -0.990 & -0.033 & negative remodeling-phosphorylation covariance summary \\
\bottomrule
\end{tabularx}
\end{table}

\begin{table}[H]
\centering\small
\caption{\textbf{Approximate future-n planning from observed effect scale.}}
\label{tab:power}
\begin{tabular}{lr}
\toprule
Mediator & Approximate n for 80\% interval-exclusion power \\
\midrule
Endothelial & 40 animals \\
Matrix/endothelial composite & 40 animals \\
Stromal/fibroblast & 60 animals \\
DCT identity & 80 animals \\
\bottomrule
\end{tabular}
\end{table}

\noindent
Interpretation: the data are consistent with remodeling-linked covariance or shared remodeling, not causal evidence. The direct flight effect remains negative after intermediate-score adjustment.

% =====================================================================
\section{External spatial reference contextualization}

\begin{figure}[H]
\centering
\includegraphics[width=\textwidth]{v11_spatial_reference_projection.pdf}
\caption{\textbf{External IRI spatial reference projection.} The OSD-462 RNA vector aligns strongly with Visium IRI repair-stage pathway programs. Spatial-niche matches are strongest in injured tubule, DCT-adjacent, fibro-inflammatory, and endothelial contexts.}
\label{fig:spatial}
\end{figure}

\begin{table}[H]
\centering\small
\caption{\textbf{Visium IRI timepoint cosines vs sham.}}
\label{tab:visiumtime}
\begin{tabular}{lrrrr}
\toprule
IRI timepoint & OSD-462 pathway & RRRM-2 pathway & OSD-462 gene & RRRM-2 gene \\
\midrule
4 h & 0.627 & 0.709 & 0.041 & 0.079 \\
12 h & 0.780 & 0.699 & 0.063 & 0.054 \\
day 2 & 0.792 & 0.667 & 0.061 & 0.090 \\
day 14 & 0.797 & 0.720 & 0.083 & 0.070 \\
week 6 & 0.801 & 0.687 & 0.028 & 0.275 \\
\bottomrule
\end{tabular}
\end{table}

\begin{table}[H]
\centering\small
\caption{\textbf{Top OSD-462 Visium niche pathway matches.}}
\label{tab:niche}
\begin{tabularx}{\textwidth}{llr}
\toprule
Timepoint & Niche & Cosine \\
\midrule
day 14 & injured tubule-enriched spots & 0.827 \\
day 14 & DCT-adjacent spots & 0.825 \\
week 6 & DCT-adjacent spots & 0.814 \\
day 14 & fibro-inflammatory repair-enriched spots & 0.807 \\
week 6 & injured tubule-enriched spots & 0.802 \\
week 6 & endothelial-enriched spots & 0.802 \\
\bottomrule
\end{tabularx}
\end{table}

\begin{figure}[H]
\centering
\includegraphics[width=0.96\textwidth]{v11_xenium_annotation_inventory.pdf}
\caption{\textbf{GSE269719 Xenium annotation inventory.} The Xenium object provides targeted-panel annotation and cellular-neighborhood context, not whole-transcriptome pathway projection.}
\label{fig:xenium}
\end{figure}

\begin{table}[H]
\centering\small
\caption{\textbf{Selected Xenium annotation counts.}}
\label{tab:xenium}
\begin{tabularx}{\textwidth}{p{0.42\textwidth}rX}
\toprule
Annotation & Count & Note \\
\midrule
Total cells & 1,374,915 & processed AnnData object loaded in backed mode \\
Genes & 299 & targeted panel, not genome-wide \\
DCT cells & 48,774 & \code{celltype\_plot} annotation \\
Endothelial cells & 129,261 & \code{celltype\_plot} annotation \\
Fibroblasts & 160,328 & \code{celltype\_plot} annotation \\
Injured proximal tubule cells & 146,097 & \code{celltype\_plot} annotation \\
Distal tubule niche cells & 77,103 & \code{CN} annotation \\
Fibro-inflammatory niche cells & 105,536 & \code{CN} annotation \\
\bottomrule
\end{tabularx}
\end{table}

\noindent
Interpretation: spatial evidence is a contextual supplement. It supports the wording ``spatially contextualized,'' not ``spatially validated.''

% =====================================================================
\section{Public perturbation-reference analyses}
\label{sec:perturbation}

\begin{figure}[H]
\centering
\includegraphics[width=\textwidth]{v11_perturbation_triangulation.pdf}
\caption{\textbf{Public perturbation-reference analyses.} (A) Low-K DCT response cosines with OSD-462, OSD-513, and RRRM-2 ISS-T RNA by gene-prior class. (B) Focused DCT/WNK transport-target gene direction heatmap. (C) Parent-protein-normalized DCT1 top-decile enrichment ORs. (D) IRI Visium DCT-transport score by timepoint and spatial context.}
\label{fig:perturbation}
\end{figure}

\begin{table}[H]
\centering\small
\caption{\textbf{Public perturbation-reference branch verdicts.}}
\label{tab:triagebranches}
\begin{tabularx}{\textwidth}{p{0.22\textwidth}p{0.30\textwidth}Xp{0.18\textwidth}}
\toprule
Mechanism branch & Reference and analysis & Headline result & Paper role \\
\midrule
Low-K/WNK reference & GSE228367 low-K DCT-enriched pseudobulk anti-alignment & Directional anti-alignment in 14-gene transport-target subset (cosine $-0.286$ OSD-462, $-0.311$ RRRM-2 ISS-T, $-0.146$ OSD-513), wide bootstrap CIs, did not meet promotion rule. & Mechanism triage, not primary upstream mechanism. \\
Parent-protein normalization & OSD-462 parent-protein-normalized phosphosite re-test & DCT1 top-decile OR 1.52 (CI 1.36--1.70, q=$1.04\times10^{-11}$); DCT2-bottom-decile also enriched; paired sample-level contrast supportive. & Robustness upgrade to the DCT-subtype-prior result; not direct biochemical occupancy. \\
Injury/repair spatial context & GSE269622 IRI Visium DCT-transport score by spatial niche & DCT-adjacent spots drop at day 14 ($-0.044$) and week 6 ($-0.034$); DCT-marker-high spots rise at the same timepoints. & Spatial prediction/context for future spaceflight kidney sections. \\
Vasopressin/cAMP & PXD001729 mpkDCT dDAVP & 60 shared single sites, 0 shared transport-target sites; key DCT/WNK targets do not map. & Secondary negative; coverage-bounded. \\
KLHL3/CUL3 turnover & OSD-462 KLHL3/CUL3/WNK coverage & 0 \gene{Klhl3} phosphosites, 1 \gene{Cul3} site, flat WNK/SPAK/NCC protein abundance; no public kidney ubiquitinomics. & Future-data rationale. \\
\bottomrule
\end{tabularx}
\end{table}

\subsection{Branch 1: GSE228367 low-K DCT anti-alignment}

\begin{table}[H]
\centering\small
\caption{\textbf{Low-K DCT response vs spaceflight RNA, 14-gene transport-target subset.}}
\label{tab:lowktargets}
\begin{tabular}{lrrrrrr}
\toprule
Spaceflight vector & $n$ genes & Cosine & 95\% CI & Spearman $\rho$ & $p$ & q \\
\midrule
OSD-462 RNA & 14 & $-0.286$ & $-0.752$ to $0.458$ & $-0.420$ & 0.135 & 0.218 \\
RRRM-2 ISS-T RNA & 14 & $-0.311$ & $-0.780$ to $0.325$ & $-0.341$ & 0.233 & 0.306 \\
OSD-513 RNA stat & 14 & $-0.146$ & $-0.768$ to $0.506$ & $-0.169$ & 0.563 & 0.695 \\
\bottomrule
\end{tabular}
\end{table}

\noindent
Interpretation: the 14-gene transport-target subset moves directionally opposite the canonical low-K DCT activation direction in all three spaceflight cohorts (10 of 14 focused DCT/WNK genes flip in OSD-462, including \gene{Slc12a3}, \gene{Wnk1}, \gene{Klhl3}, \gene{Cul3}), but none of the comparisons clears the pre-specified primary threshold (cosine $<-0.3$ with CI excluding zero). Promotion rule not met; treated as mechanism triage.

\subsection{Branch 2: parent-protein-normalized DCT-prior enrichment}

\begin{table}[H]
\centering\small
\caption{\textbf{Parent-protein-normalized phosphosite enrichment among DCT-prior parent genes.} Parent-normalized effect $=$ phosphosite flight effect minus parent-protein flight effect. This is a robustness summary, not direct biochemical occupancy.}
\label{tab:occupancy}
\begin{tabularx}{\textwidth}{p{0.26\textwidth}p{0.18\textwidth}rrrX}
\toprule
Analysis & DCT-prior subset & OR & 95\% CI & q & Read \\
\midrule
parent-normalized $p<0.05$ & DCT1 top decile & 1.52 & 1.36--1.70 & $1.04\times10^{-11}$ & robustness upgrade holds \\
parent-normalized $p<0.05$ & DCT2 bottom decile & 1.32 & 1.18--1.48 & $4.27\times10^{-6}$ & comparator also enriched \\
parent-normalized strict $q<0.10$ & DCT1 top decile & 1.61 & 1.42--1.82 & $1.04\times10^{-11}$ & stricter set, OR rises \\
parent-normalized strict $q<0.10$ & DCT2 bottom decile & 1.28 & 1.13--1.46 & $3.01\times10^{-4}$ & comparator also enriched \\
exclude anchor genes & DCT1 top decile & 1.51 & 1.34--1.69 & $2.73\times10^{-11}$ & survives anchor exclusion \\
exclude NCC sites & DCT1 top decile & 1.51 & 1.35--1.69 & $2.73\times10^{-11}$ & survives NCC exclusion \\
one row per gene & DCT1 top decile & 1.50 & 1.19--1.90 & 0.00101 & survives parent-gene representative-row sensitivity \\
single-position row per gene & DCT1 top decile & 1.48 & 1.17--1.88 & 0.00147 & stricter representative-row sensitivity passes \\
paired sample-level contrast & DCT1 top decile & 1.56 & 1.37--1.77 & $1.90\times10^{-10}$ & matched animal-level contrast supportive \\
paired sample-level contrast & DCT2 bottom decile & 1.54 & 1.36--1.75 & $2.22\times10^{-10}$ & matched animal-level comparator supportive \\
\bottomrule
\end{tabularx}
\end{table}

\begin{table}[H]
\centering\small
\caption{\textbf{Selected target-site parent-protein-normalized examples.}}
\label{tab:occupancytarget}
\begin{tabular}{llrrr}
\toprule
Gene & Site & Raw phospho effect & Parent protein effect & Parent-normalized effect \\
\midrule
\gene{Slc12a3}/NCC & p53 & $-0.851$ & $+0.089$ & $-0.940$ \\
\gene{Slc12a3}/NCC & p65 & $-0.790$ & $+0.089$ & $-0.879$ \\
\gene{Slc12a3}/NCC & p68 & $-0.794$ & $+0.089$ & $-0.883$ \\
\gene{Slc12a3}/NCC & p89 & $-0.930$ & $+0.089$ & $-1.019$ \\
\gene{Stk39}/SPAK & p366 & $-0.520$ & $+0.081$ & $-0.601$ \\
\gene{Stk39}/SPAK & p383 & $-0.793$ & $+0.081$ & $-0.874$ \\
\bottomrule
\end{tabular}
\end{table}

\noindent
Interpretation: per-site parent-protein normalization sharpens the suppressed-site effect because parent NCC and SPAK protein effects are slightly positive. The DCT-prior enrichment is not eliminated by parent protein abundance in this sensitivity; the calculation remains a whole-kidney parent-normalized effect rather than direct phosphosite occupancy or biochemical stoichiometry.

\subsection{Branch 3: IRI Visium DCT-transport spatial context}

\begin{table}[H]
\centering\small
\caption{\textbf{Visium DCT-transport score versus sham, by spatial group and timepoint.} Spot-level descriptive two-sample $t$; no animal-level replication.}
\label{tab:dctxport}
\begin{tabular}{lllrrr}
\toprule
Spatial group & Timepoint & Mean diff vs sham & $t$ & $p$ & $n$ spots vs sham \\
\midrule
DCT-adjacent spots & day 14 & $-0.044$ & $-2.81$ & 0.0050 & 932 vs 1{,}134 \\
DCT-adjacent spots & week 6 & $-0.034$ & $-2.36$ & 0.0182 & 1{,}574 vs 1{,}134 \\
DCT-marker-high spots (top quartile) & day 14 & $+0.199$ & $+8.02$ & $2.25\times10^{-15}$ & 693 vs 715 \\
DCT-marker-high spots (top quartile) & week 6 & $+0.132$ & $+5.98$ & $2.70\times10^{-9}$ & 979 vs 715 \\
\bottomrule
\end{tabular}
\end{table}

\noindent
Interpretation: DCT-adjacent neighborhoods carry a DCT-transport drop in late IRI repair while DCT-marker-high spots retain or even increase the DCT-transport program. This is the spatial prediction that a future spaceflight kidney spatial experiment would test directly.

\subsection{Branch 4: PXD001729 dDAVP/mpkDCT}

\begin{table}[H]
\centering\small
\caption{\textbf{PXD001729 mpkDCT dDAVP reference vs OSD-462.}}
\label{tab:pxdshared}
\begin{tabular}{lrr}
\toprule
Quantity & Value & Note \\
\midrule
Shared single phosphosites & 60 & cross-platform/cross-species mapping \\
Shared changed transport-target sites & 0 & \gene{Slc12a3}, \gene{Stk39}, \gene{Oxsr1}, \gene{Wnk1}, \gene{Wnk4}, \gene{Klhl3}, \gene{Cul3} not shared \\
All-shared-site cosine & $-0.026$ & 95\% bootstrap CI $-0.310$ to $0.287$ \\
\bottomrule
\end{tabular}
\end{table}

\noindent
Interpretation: the cultured DCT-lineage dDAVP reference is coverage-bounded for the OSD-462 transport-target lesion and cannot test a vasopressin/cAMP anti-alignment hypothesis. Used as plausibility/coverage boundary only.

\subsection{Branch 5: KLHL3/CUL3/WNK turnover}

\begin{table}[H]
\centering\small
\caption{\textbf{KLHL3/CUL3/WNK coverage in OSD-462 and PXD001729.}}
\label{tab:klhlcoverage}
\begin{tabular}{lrrrl}
\toprule
Gene & OSD-462 phosphosites & PXD phosphosites & OSD-462 protein effect & Read \\
\midrule
\gene{Klhl3} & 0 & 0 & $+0.013$ & no phosphosite evidence \\
\gene{Cul3} & 1 & 0 & $-0.001$ & no turnover signal \\
\gene{Wnk1} & 12 & 12 & $-0.024$ & phospho observed, protein flat \\
\gene{Wnk4} & 8 & 1 & $+0.087$ & phospho observed, protein not depleted \\
\gene{Stk39}/SPAK & 5 & 3 & $+0.081$ & phospho suppressed despite flat protein \\
\gene{Slc12a3}/NCC & 16 & 0 & $+0.089$ & phospho suppressed with protein flat/up \\
\bottomrule
\end{tabular}
\end{table}

\noindent
Interpretation: public data cannot distinguish KLHL3/CUL3-mediated WNK turnover from ionic/osmotic suppression. The next experiment requires DCT-enriched or targeted ubiquitinomics and direct WNK protein turnover assays.

% =====================================================================
\section{Complete artifact index}

\begin{longtable}{p{0.30\textwidth}p{0.64\textwidth}}
\caption{\textbf{V11 generated result artifacts indexed by source path.}}\\
\toprule
Category & Source artifact(s) \\
\midrule
\endfirsthead
\toprule
Category & Source artifact(s) \\
\midrule
\endhead
\bottomrule
\endfoot
Run report & \path{docs/v11_execution_research_plan.md}; \path{docs/v11_execution_results.md} \\
Headline-number index & \path{latex_paper/v11_headline_number_index.tsv}; maps reported manuscript numbers to source artifacts, filters, and columns. \\
V11 source modules & \path{src/v11/build_gse228367_dct_prior.R}; \path{src/v11/core_analysis.py}; \path{src/v11/h2_composition_aware_phospho.py}; \path{src/v11/spatial_reference_projection.py}; \path{src/v11/publication_figures.py}; \path{src/v11/perturbation_gse228367_lowk.py}; \path{src/v11/h2_occupancy_normalized_phospho.py}; \path{src/v11/spatial_dct_transport_check.py}; \path{src/v11/perturbation_triage_summary.py} \\
Runner integration & \path{src/run_all_phases.py} phase 11; \code{--v11}; \code{--v11-only}; \code{--v11-site-scope}; \code{--v11-skip-spatial}; \code{--v11-skip-visium}; \code{--v11-skip-xenium}; \code{--v11-skip-figures} \\
Baseline lock & \path{data/results/run_20260526_v11_dct1_phospho_mediation/baseline/v11_baseline_lock_summary.tsv}; \path{baseline/v11_baseline_input_manifest.json}; \path{baseline/osd462_targeted_ksea_substrates.tsv}; \path{baseline/rna_recurrence_gene_set_members.tsv}; \path{baseline/osd462_rna_recurrence_leave_one_family.tsv}; \path{baseline/osd462_rna_recurrence_paired_pathway_bootstrap.tsv}; \path{baseline/osd462_tmt_channel_qc.tsv}; \path{baseline/osd462_tmt_missingness_by_condition.tsv} \\
DCT reference prior & \path{dct_prior/gse228367_gene_stats_overall.tsv}; \path{dct_prior/gse228367_gene_stats_by_rep.tsv}; \path{dct_prior/gse228367_dct1_vs_dct2_de.tsv}; \path{dct_prior/gse228367_dct_prior_summary.json}; \path{dct_prior/gse228367_dct_prior_score_sensitivity.tsv}; \path{dct_prior/gse228367_anchor_gene_prior_positions.tsv}; \path{dct_prior/gse228367_prior_bin_stability.tsv} \\
External QC & \path{external_qc/gse228367_object_inventory.tsv}; \path{external_qc/gse228367_marker_qc.tsv}; \path{external_qc/pxd001729_table_inventory.tsv}; \path{external_qc/pxd001729_phosphosite_effects.tsv} \\
DCT-subtype-prior enrichment & \path{h2_enrichment/h2_dct1_sensitivity_summary.tsv}; \path{h2_enrichment/h2_dct1_parent_gene_level_summary.tsv}; \path{h2_enrichment/h2_dct1_site_count_stratified_permutation.tsv}; \path{h2_enrichment/h2_dct_extreme_bin_cluster_bootstrap.tsv}; \path{h2_enrichment/h2_dct_matched_parent_gene_permutation.tsv}; \path{h2_enrichment/h2_dct_threshold_free_summary.tsv}; \path{h2_enrichment/h2_dct_percentile_threshold_sensitivity.tsv}; \path{h2_enrichment/h2_dct1_enrichment_verdict.json} \\
PXD001729 checks & \path{h2_pxd/h2_pxd001729_ddavp_antialignment_summary.tsv}; \path{h2_pxd/h2_pxd001729_shared_sites.tsv}; \path{h2_pxd/h2_pxd001729_verdict.json} \\
KLHL3/CUL3 checks & \path{h2_klhl3/h2_klhl3_cul3_effects.tsv}; \path{h2_klhl3/h2_klhl3_cul3_verdict.json} \\
Composition-aware models & \path{h2_composition_adjusted/h2_composition_adjusted_suppression_enrichment_single.tsv}; \path{h2_composition_adjusted/h2_composition_effect_level_dct1_ladder_single.tsv}; \path{h2_composition_adjusted/h2_composition_site_fixed_interaction_ladder_single.tsv}; \path{h2_composition_adjusted/h2_composition_fixed_set_ladder_single.tsv}; \path{h2_composition_adjusted/h2_composition_covariate_diagnostics_single.tsv}; \path{h2_composition_adjusted/h2_composition_aware_verdict_single.json}; \path{h2_composition_adjusted/h2_composition_manifest_single.json} \\
Exploratory covariance decomposition & \path{h3_mediation/h3_mediation_model_summary.tsv}; \path{h3_mediation/h3_mediation_posterior_draws.tsv.gz}; \path{h3_mediation/h3_mediation_power_simulation.tsv}; \path{h3_mediation/h3_mediation_verdict.json} \\
Spatial reference & \path{spatial_reference/spatial_reference_projection_verdict.json}; \path{spatial_reference/visium_timepoint_gene_cosines.tsv}; \path{spatial_reference/visium_timepoint_pathway_cosines.tsv}; \path{spatial_reference/visium_niche_cosines.tsv}; \path{spatial_reference/visium_spot_niche_counts.tsv}; \path{spatial_reference/xenium_annotation_inventory.tsv}; \path{spatial_reference/visium_dct_transport_by_niche.tsv}; \path{spatial_reference/visium_dct_transport_timepoint_summary.tsv}; \path{spatial_reference/visium_dct_transport_spot_scores.tsv.gz}; \path{spatial_reference/visium_dct_transport_verdict.json} \\
Perturbation triangulation & \path{perturbation/lowk_dct_pseudobulk_effects.tsv}; \path{perturbation/lowk_dct_alignment_summary.tsv}; \path{perturbation/lowk_target_gene_table.tsv}; \path{perturbation/lowk_dct1_dct2_specificity.tsv}; \path{perturbation/lowk_alignment_verdict.json}; \path{perturbation/perturbation_triage_summary.tsv}; \path{perturbation/perturbation_triage_verdict.json} \\
Parent-protein-normalized re-test & \path{h2_occupancy/h2_occupancy_site_effects.tsv}; \path{h2_occupancy/h2_occupancy_dct1_enrichment.tsv}; \path{h2_occupancy/h2_occupancy_target_sites.tsv}; \path{h2_occupancy/h2_parent_normalized_paired_site_effects.tsv.gz}; \path{h2_occupancy/h2_parent_normalized_paired_enrichment.tsv}; \path{h2_occupancy/h2_occupancy_verdict.json} \\
V11 figures & \path{figures/v11/v11_dct1_parent_gene_enrichment.pdf}; \path{figures/v11/v11_distal_nephron_prior_enrichment.pdf}; \path{figures/v11/v11_parent_protein_composition_sensitivity.pdf}; \path{figures/v11/v11_exploratory_mediation_forest.pdf}; \path{figures/v11/v11_exploratory_covariance_decomposition.pdf}; \path{figures/v11/v11_spatial_reference_projection.pdf}; \path{figures/v11/v11_xenium_annotation_inventory.pdf}; \path{figures/v11/v11_perturbation_triangulation.pdf} \\
Downloaded spatial data & \path{data/external/spatial_reference/GSE269622_Visium/*.tar.gz}; \path{data/external/spatial_reference/GSE269719_Xenium/Xenium.h5ad} \\
\end{longtable}

% =====================================================================
\section{Verification notes}

\begin{itemize}
  \item The v11 figure module generated PNG/PDF figure sets under \path{figures/v11/}, including the distal-nephron prior enrichment alias, parent-protein composition sensitivity, exploratory covariance-decomposition alias, spatial-reference projection, Xenium annotation inventory, and perturbation triangulation.
  \item The direct v11 figure generator and integrated \path{src.visualization.publication_figures} path both detect and regenerate v11 figures.
  \item Python syntax checks passed for the v11 core, composition-aware, spatial-reference, figure, perturbation low-K, parent-protein-normalized phosphosite, spatial DCT-transport, and triage modules, plus the runner and publication-figure integration.
  \item The runner dry-run passed for \code{python -m src.run_all_phases --v11-only --dry-run --run-id run\_20260526\_v11\_dryrun --v11-skip-spatial} and for the perturbation-triangulation modules.
  \item The GSE269719 Xenium MD5 matched the figshare record: \code{e6f5728ab7e47b499149f63b54e892c7}.
  \item The spatial scripts completed after loading six GSE269622 Visium archives and the processed GSE269719 Xenium AnnData object; the DCT-transport spatial check used the same Visium archives.
  \item The perturbation low-K analysis loaded six GSE228367 filtered 10x matrices (NK1--3 and KD1--3) and emitted alignment and target-gene tables under \path{perturbation/}.
\end{itemize}

\vspace{1em}
\noindent\itshape
Curated results compendium for Manuscript v11. This document should travel with \path{manuscript_v11.tex} and \path{manuscript_v11.pdf}. The manuscript result is a distal-nephron subtype-prior phosphosite enrichment model embedded in recurrent kidney remodeling.

\end{document}
